\section{Introduction}

When a Qwen2.5-7B-Instruct agent attempts to boil water in ScienceWorld \cite{scienceworld-2022}, it must track temperature thresholds, container states, and heat source interactions across dozens of steps. Internal LLM reasoning hallucinates these state transitions---predicting water has boiled when the stove was never activated---triggering agent death with a $-100$ penalty. Deterministic external computation (environment simulators, code interpreters) eliminates such hallucinations by executing exact state transitions. The natural question: when should an agent route computation externally versus reasoning internally?

Existing routing approaches \cite{toolformer-2024,art-2023,chameleon-2024} frame tool selection as a capability or efficiency decision---matching subtasks to tools that can solve them. This framing ignores a potentially dominant signal: routing to \emph{avoid catastrophic failure}. Understanding what routing actually learns requires counterfactual analysis, but the standard per-step oracle approach degenerates under sparse rewards. On ScienceWorld, where non-zero reward arrives only at episode termination, 97.4\% of steps produce tied intermediate scores under both strategies across 245 episodes---no step-level signal exists to train a router (\S\ref{sec:collection}).

We propose \emph{episode-level counterfactual analysis}: run complete episodes under pure-internal (ReAct~\cite{react-2023}) and pure-deterministic strategies, then compare episode outcomes to generate routing labels (\S\ref{sec:collection}). Across 11 ScienceWorld task types (245 episodes), the full deterministic advantage is $\Delta_{\text{full}} = +29.60$ points. Excluding death episodes, this gap collapses to $\Delta_{\text{nd}} = +2.84$---a Death Avoidance Factor (DAF; \S\ref{sec:daf}) of 90.41\%. The per-task correlation between death rate differential and performance gap is $r = 0.9698$ ($p = 8.12 \times 10^{-7}$), confirming death avoidance as the primary mechanism (\S\ref{sec:death-decomp}). Death rates differ sharply: 42.9\% for internal reasoning versus 16.3\% for deterministic computation. Of 245 episodes, 86 prefer deterministic, 9 prefer internal, and 150 show no preference---96.3\% of oracle decisions match always-deterministic (\S\ref{sec:baselines}).

Our contributions are twofold: (1)~an episode-level counterfactual methodology with death-aware labeling and the DAF metric, overcoming per-step oracle degeneracy (\S\ref{sec:collection}, \S\ref{sec:daf}); and (2)~the empirical finding that DAF = 90.41\% across 245 episodes, confirmed by router training as diagnostic probe (\S\ref{sec:router-diagnostic}).
